%*************************************************************************
% Grundlagen
%*************************************************************************
\chapter{Grundlagen}
\label{ch:grundlagen}

\section{Einleitung}
\label{sec:grundlagen-einleitung}

Dieses Kapitel bietet einen kurzen Einstieg in die zentralen Konzepte, auf denen
die weitere Arbeit aufbaut. Zu Beginn werden die von \textcite{dung:1995:af}
eingeführten \emph{Argumentationssysteme} vorgestellt. 

Anschließend werden verschiedene Semantiken betrachtet, mit denen sich die
Akzeptanz von Argumenten bestimmen lässt. Neben der extensionsbasierten Sicht
wird insbesondere die labellingbasierte Darstellung eingeführt, da sie im
weiteren Verlauf der Arbeit die Grundlage für die logische Kodierung der
Unabhängigkeitsbegriffe bildet.

Darauf aufbauend wird das Konzept der \emph{bedingten Unabhängigkeit}
vorgestellt. In den späteren Kapiteln wird dieser Begriff auf
Argumentationssysteme übertragen.

Abschließend werden \emph{Quantifizierte Boolesche Formeln} eingeführt. Sie
bilden das formale Werkzeug, mit dem die Unabhängigkeitsbegriffe für 
verschiedene Argumentationssemantiken kodiert werden können.

\section{Argumentationssysteme}
\label{sec:af}

\emph{Argumentationssysteme} (engl. \emph{Argumentation Framework}, AF)
sind ein zentrales Konzept der formalen Argumentationstheorie. Sie wurden von
Dung~\cite{dung:1995:af} eingeführt, um argumentative Strukturen abstrakt zu
modellieren. Der Grundgedanke besteht darin, Argumente als abstrakte Einheiten
zu betrachten, die durch eine Angriffsrelation miteinander verbunden sind.
Die Akzeptanz eines Arguments hängt damit nicht von seiner inneren Struktur ab,
sondern davon, wie es zu anderen Argumenten im System in Beziehung steht.

Die grundlegenden Definitionen dieses Abschnitts orientieren sich, sofern nicht
anders angegeben, an \textcite{handbook__of_formal_argumentation__vol__3:2024}.
Die labellingbasierte Perspektive sowie die dafür verwendete Notation folgen in
Anlehnung an \textcite{bluemel:2025:independence-aba}.

\begin{Definition}
    \label{def:af}
    Ein \emph{Argumentationssystem} (AF) ist ein Paar \(F=(A,R)\), wobei
    \(A\) eine Menge von Argumenten und
    \(R\subseteq A\times A\) eine binäre Angriffsrelation auf \(A\) ist.
    Für \((a,b)\in R\) sagen wir, dass \(a\) das Argument \(b\) \emph{angreift}.
    Ferner schreiben wir $A_F$ und $R_F$, um die Zugehörigkeit zum AF F auszudrücken.
    Ist dies im Kontext klar, verzichten wir auf diese Kennzeichnung. 
\end{Definition}

Grundsätzlich interessiert uns insbesondere, welche Argumente andere Argumente
angreifen und welche Argumente durch eine Menge von Argumenten verteidigt
werden. Dazu führen wir zunächst folgende Notation ein.

\begin{Definition}
    \label{attack-object-subject}
    Sei \(F=(A,R)\) ein AF, sei \(S\subseteq A\) und sei \(a\in A\). Wir
    definieren:
    \[
        a^+ = \{b\in A \mid (a,b)\in R\},~
        S^+ = \{b\in A \mid \exists a\in S:(a,b)\in R\},
    \]
    \[
        a^- = \{b\in A \mid (b,a)\in R\},~
        S^- = \{b\in A \mid \exists a\in S:(b,a)\in R\}.
      \]
\end{Definition}

Dabei bezeichnet \(a^+\) die Menge aller Argumente, die von \(a\) angegriffen
werden, und \(a^-\) die Menge aller Angreifer von \(a\). Die Mengen \(S^+\)
und \(S^-\) übertragen diese Notation auf Mengen von Argumenten.

Mit dieser Notation lässt sich das Konzept der Verteidigung wie folgt definieren.
\begin{Definition}
    \label{def:defense}
    Sei \(F=(A,R)\) ein AF. Ein Argument \(a\in A\) heißt
    \emph{verteidigt} durch eine Menge \(S\subseteq A\), genau dann wenn
    für jedes Argument \(b\in a^-\) ein Argument \(c\in S\) existiert, sodass
      \(b\in c^+\) gilt.
\end{Definition}

Insbesondere kann auch die leere Menge ein Argument verteidigen. Dies ist genau
dann der Fall, wenn ein Argument keinen Angreifer besitzt, also \(a^-=\{\}\)
gilt.

Die Frage, welche Argumente in einem Argumentationsszenario als überzeugend gelten sollen, 
wird mit dem Konzept der \emph{Argumentationssemantiken} beantwortet.
Eine jede Semantik lässt sich dabei aus einer \emph{labelling}- und einer \emph{extensions}-basierten 
Perspektive formulieren. 

Der naheliegende Ansatz sich aus der Argumentmenge eine Auswahl an Argumenten herauszusuchen, die gemeinsam das
Argumentationsszenario \enquote{überleben}, führt zum \emph{extensions}-basierten Ansatz. 
Eine \emph{Extension} ist dann genau die ausgewählte Argumentmenge. Je nach Semantik werden unterschiedliche 
und auch mehrere Kriterien gleichzeitig verwandt, um eine Menge von Argumenten auszuwählen. 
\begin{Definition}
      \label{def:extension}
      Eine \emph{Extension} ist eine Menge \(E\subseteq A\), die bezüglich einer
      gegebenen Semantik bestimmte Akzeptanzbedingungen erfüllt.
\end{Definition}
\begin{Definition}
      \label{def:semantik}
      Eine \emph{extensions}-basierte Semantik $\sigma$ ist eine Funktion, die jedem AF \(F\) eine Menge an Extensionen zuordnet. 
      Es gilt: $\sigma(F) \subseteq 2^{A}$
\end{Definition}

Der \emph{labellingbasierte Ansatz} bewertet jedes Argument einzeln durch ein
Label. Im Rahmen dieser Arbeit verwenden wir die drei Labels
\(\mathit{in}\), \(\mathit{out}\) und \(\mathit{und}\). Dabei steht
\(\mathit{in}\) intuitiv für akzeptierte Argumente und \(\mathit{out}\) für
verworfene Argumente. Das Label \(\mathit{und}\) beschreibt Argumente, deren
Status unentschieden bleibt.

Eine labellingbasierte Semantik ordnet einem AF nicht nur eine einzelne
Bewertung, sondern eine Menge zulässiger Labellings zu. Ein Labelling ist dabei
eine Funktion, die jedem Argument genau eines der drei Labels zuweist.
\begin{Definition}
      \label{def:labelling}
      Sei \(F=(A,R)\) ein AF. Ein \emph{Labelling} von \(F\) ist eine Funktion
      \(\lambda : A \to \{\mathit{in},\mathit{out},\mathit{und}\},\)
      die jedem Argument genau eines der drei Labels
      \(\mathit{in}\), \(\mathit{out}\) oder \(\mathit{und}\) zuordnet.
\end{Definition}

Je nach Semantik gelten unterschiedliche Bedingungen dafür, welche Labellings
als zulässig betrachtet werden. Die Menge aller Labellings, die eine Semantik
\(\sigma\) einem Argumentationssystem \(F\) zuordnet, wird im Folgenden mit
\(\Lambda_\sigma(F)\) bezeichnet.
\begin{Definition}
      \label{def:labelling-set-of-all}
      Sei \(F=(A,R)\) ein AF und sei \(\sigma\) eine labellingbasierte
      Semantik. Für die Menge aller \(\sigma\)-Labellings von \(F\) schreiben
      wir \(\Lambda_\sigma(F).\)
\end{Definition}

Für die spätere Definition bedingter Unabhängigkeit ist es wichtig, Labellings
nicht nur auf der gesamten Argumentmenge \(A\), sondern auch auf Teilmengen von
Argumenten betrachten zu können. Dazu verwenden wir die Einschränkung eines
Labellings auf eine Teilmenge \(X\subseteq A\).
\begin{Definition}
      \label{def:labelling-rest}
      Sei \(F=(A,R)\) ein AF, sei \(X\subseteq A\) und sei \(\lambda\) ein
      Labelling von \(F\). Mit \(\lambda|_X\) bezeichnen wir die Einschränkung
      von \(\lambda\) auf \(X\), also die Funktion \(\lambda|_X : X \to \{\mathit{in},\mathit{out},\mathit{und}\}\)
      mit \(\lambda|_X(a)=\lambda(a)~ \text{für alle } a\in X.\)

\end{Definition}
Häufig interessiert nicht nur, welches Label ein einzelnes Argument besitzt,
sondern auch die Menge aller Argumente mit einem bestimmten Label. Dafür führen
wir die folgende Notation ein.

\begin{Definition}
      \label{def:labelling-set}
      Sei \(F=(A,R)\) ein AF, sei
      \(L=\{\mathit{in},\mathit{out},\mathit{und}\}\)
      die Menge der Labels und sei \(\lambda\) ein Labelling von \(F\).
      Für ein Label \(lab\in L\) definieren wir
      \(lab(\lambda)=\{a\in A\mid \lambda(a)=lab\}.\)
\end{Definition}
\newpage
\section{Semantiken}
\label{sec:af-semantiken}



Basis aller Semantiken ist die \emph{Konfliktfreiheit}. Eine Menge an Argumenten soll ein Argumentationsszenario nur "überleben", 
wenn es in sich selbst stimmig ist \cite{handbook__of_formal_argumentation__vol__3:2024}.
\begin{Definition}
    \label{def:conflict-free}
    Sei $F = (A, R)$ ein AF. Eine Menge $S \subseteq A$ heißt \emph{konfliktfrei}, genau dann wenn es keine Argumente $a, b \in S$ gibt, 
    so dass $a \rightarrow b$ gilt.
\end{Definition}
Insbesondere ist damit ein Argument, das sich selbst angreift, nicht \emph{konfliktfrei}.

\begin{Definition}
      \label{range}
      Sei $F = (A, R)$ ein AF und $S \subseteq A$. Die \emph{Reichweite} von $S$ ist definiert als $S \cup S^{+}$.
\end{Definition}

\begin{Definition}
      \label{semantic-extensions}
      Sei $F = (A, R)$ ein AF, $S \subseteq A$ und $S$ \emph{konfliktfrei}. $S$ heißt:
\begin{itemize}
      \item \emph{admissible}, genau dann wenn jedes Argument in $S$ durch $S$ verteidigt wird.
      \item \emph{complete}, genau dann wenn jedes Argument, das durch $S$ verteidigt wird, in $S$ enthalten ist.
      \item \emph{preferred}, genau dann wenn $S$ inhaltsmaximal unter den admissible Mengen ist.
      \item \emph{stable}, genau dann wenn $S^{+} = A \setminus S$ gilt.
\end{itemize}
\end{Definition}


Eine weitere Möglichkeit, die Akzeptanz von Argumenten zu bestimmen, 
ist der labelling-basierte Ansatz, 
der von Caminada~\cite{caminada:2009:logical-account} eingeführt wurde.


Wir verwenden $\text{in}(\mathcal{L})$ für $\{a \mid \mathcal{L}(a) = \text{in}\}$, 
$\text{out}(\mathcal{L})$ für $\{a \mid \mathcal{L}(a) = \text{out}\}$ 
und $\text{und}(\mathcal{L})$ für $\{a \mid \mathcal{L}(a) = \text{und}\}$.

\begin{Definition}
      \label{def:complete-labelling}
      Sei $F = (A, R)$ ein AF und $\mathcal{L} : A \rightarrow \{\text{in}, \text{out}, \text{und}\}$ ein $labelling$
      von $F$. $\mathcal{L}$ heißt \emph{complete}, wenn für jedes Argument $a \in A$ gilt:
      \begin{itemize}
            \item $\mathcal{L}(a) = \text{in}$ gdw. für jedes $b \in A$ mit $(b, a) \in R$ $\mathcal{L}(b) = \text{out}$ gilt.
            \item $\mathcal{L}(a) = \text{out}$ gdw. es ein $b \in A$ mit $(b, a) \in R$ gibt, für das $\mathcal{L}(b) = \text{in}$ gilt.
      \end{itemize}
\end{Definition}

\section{Conditional Independence}
\label{sec:conditional-independence}
      Der Begriff der \emph{conditional independence} ist ein zentrales Konzept in der Wahrscheinlichkeitstheorie.
      Er beschreibt die Unabhängigkeit von zwei Zufallsvariablen unter der Bedingung, 
      dass eine andere dritte Variable einen bestimmten Wert annimmt.

      Die semi-graphoiden Axiome, nach Pearl~\cite{pearl:2010}, sind eine Reihe von Eigenschaften, 
      die die konditionale Unabhängigkeitsrelation erfüllt.\par Sie lauten wie folgt:
      \begin{itemize}
            \item Symmetrie: $A \perp B \mid C \implies B \perp A \mid C$
            \item Dekomposition: $A \perp B \cup D \mid C \implies A \perp B \mid C$ 
            \item Schwache Vereinigung: $A \perp B \cup D \mid C \implies A \perp B \mid C \cup D$ 
            \item Kontraktion: $A \perp B \mid C \land A \perp D \mid C \cup B \implies A \perp B \cup D \mid C$
      \end{itemize}


\section{Quantified Boolean Formulas}
\label{sec:qbf}
      \emph{Quantified boolean formulas} (QBF) sind eine Erweiterung der klassischen Aussagenlogik, 
      bei der Quantoren über boolesche Variablen verwendet werden.
      \begin{Definition}
      \label{def:qbf}
      Sei $A$ ein boolescher Ausdruck. Ein \emph{QBF} ist 
      \end{Definition}

      Für eine endliche Menge aussagenlogischer Variablen
      \(V=\{v_1,\dots,v_n\}\) schreiben wir
      \[
      \forall V\,\varphi
      \]
      als Abkürzung für
      \[
      \forall v_1\dots \forall v_n\,\varphi.
      \]
      Analog verwenden wir \(\exists V\).

      Ebenso nutzen wir für zwei endliche Mengen aussagenlogischer Variablen $X=\{x_1,\dots,x_n\}, Y=\{y_1,\ldots,y_n\}$
      \[
      \forall X,Y \varphi
      \quad \text{als Abkürzung für}\quad
      \forall X \forall Y \varphi
      \]